\section{Introducción Teórica}
\label{sec:intro}

El modelo de red neuronal de Hopfield (1982) es una red recurrente y completamente conectada que funciona como una memoria autoasociativa. El sistema está compuesto por $N$ neuronas binarias (o bipolares), cuyos estados se definen como $S_i \in \{-1, 1\}$. 

\subsection{Regla de Aprendizaje y Almacenamiento}
Para almacenar un conjunto de $P$ patrones (o memorias) definidos como $\xi^\mu$ (con $\mu = 1, \dots, P$), se ajustan los pesos sinápticos $W_{ij}$ entre las neuronas utilizando una generalización de la regla de Hebb. Según \cite{hertz1991}, el peso de la conexión entre la neurona $j$ y la neurona $i$ está dado por:

\begin{equation}
	W_{ij} = \frac{1}{N} \sum_{\mu=1}^{P} \xi_i^\mu \xi_j^\mu \quad \text{con } W_{ii} = 0
	\label{eq:hebb}
\end{equation}

Esta regla genera una matriz de pesos estrictamente simétrica ($W_{ij} = W_{ji}$) sin autoconexiones, lo cual es fundamental para garantizar la estabilidad del sistema.

\subsection{Dinámica de Actualización}
El estado de la red evoluciona en el tiempo dependiendo del campo local $h_i$ que recibe cada neurona, definido como la suma ponderada de los estados de las demás neuronas conectadas a ella:

\begin{equation}
	h_i = \sum_{j \neq i} W_{ij} S_j
	\label{eq:campo_local}
\end{equation}

La regla de actualización determina que la neurona toma el signo de su campo local: $S_i = \text{sgn}(h_i)$. Esta actualización puede realizarse de forma sincrónica (todas las neuronas al mismo tiempo) o de forma asincrónica (una neurona elegida al azar por vez).

\subsection{Función de Energía y Atractores}
Una de las propiedades más importantes de la red de Hopfield es la existencia de una función de energía (o función de Lyapunov). Para una matriz de pesos simétrica, la energía global del sistema se define como:

\begin{equation}
	E = -\frac{1}{2} \sum_{i \neq j} W_{ij} S_i S_j
	\label{eq:energia}
\end{equation}

Bajo actualización asincrónica, cualquier cambio en el estado de una neurona asegura que la energía total del sistema disminuya o se mantenga igual ($\Delta E \leq 0$). Esto garantiza que la red siempre converja hacia un mínimo local de energía (atractor). Los patrones memorizados $\xi^\mu$ actúan como estos atractores. Sin embargo, el sistema también crea mínimos no deseados, conocidos como \textit{estados espurios}.

\subsection{Estados Espurios}
\label{subsec:espurios}

En una red de Hopfield, se busca que los patrones almacenados $\xi^\mu$ actúen como mínimos absolutos en el paisaje de energía. Sin embargo, la dinámica del sistema y la regla de aprendizaje de Hebb también generan mínimos locales no deseados, conocidos como \textit{estados espurios}. Si la red converge a uno de estos valles, el sistema recupera una memoria estable que no pertenece al conjunto original de patrones enseñados.

Según la teoría desarrollada por \cite{hertz1991}, estos estados aparecen inevitablemente y se clasifican en tres tipos principales:
\begin{itemize}
	\item \textbf{Estados inversos:} Dado que la función de energía es perfectamente simétrica respecto al signo de los estados ($E(S) = E(-S)$), por cada patrón original memorizado, la red también adopta como atractor a su inverso o ``negativo'' exacto ($-\xi^\mu$).
	\item \textbf{Estados de mezcla (\textit{mixture states}):} Se originan por la interferencia entre memorias. Son atractores espurios formados por la combinación lineal de un número impar de patrones almacenados (por ejemplo, una mezcla de tres patrones distintos).

\end{itemize}

\subsection{Capacidad}
\label{subsec:capacidad}
La capacidad de la red (el número máximo de patrones $P$ que se pueden recuperar con un error aceptable) depende de la interferencia o \textit{crosstalk} entre los distintos patrones memorizados. La capacidad normalizada se define como $\alpha = P / N$. 

Cuando los patrones son pseudo-aleatorios y están descorrelacionados, las interferencias se cancelan estadísticamente y la capacidad máxima teórica es de aproximadamente $\alpha \approx 0.138$. Sin embargo, si los patrones presentan correlación estructural, el término de \textit{crosstalk} aumenta drásticamente, deformando el paisaje de energía, reduciendo la profundidad de los atractores y disminuyendo la capacidad efectiva de la red. La eliminación de sinapsis (daño estructural) produce un efecto análogo, allanando los pozos de energía y facilitando que la red decante en estados espurios frente al ruido.